\section{Discussion}\label{sec:discussion}

\subsection{Who votes beats what the case is about}

The M3a null result is a negative finding with a positive reading.
The failure of three model families to convert issue areas,
lower-court dispositions, terms, and argument lengths into votes
tells us that the mapping from case facts to outcomes is either
unstable across time or not expressible in these coarse codes. The
ablation sharpens this: keeping issue area actively hurt
generalization, which is the signature of a signal that shifted
between the training and evaluation windows---the mix and meaning
of ``economic activity'' cases in OT2015 is not the mix of OT2022.
Justice identity, by contrast, travels perfectly across the split,
because it is a property of the voter, not of the era.

This is a modern, pre-registered echo of the attitudinal model's
central claim~\citep{segal2002}, obtained with none of its
theoretical apparatus: one revealed-ideology bit per justice
carries more predictive weight than the entire structured case
file. It also explains, mechanically, why the earlier literature's
headlines sit above our bar. \citet{katz2017} report 71.9\% at the
vote level---but over 1816--2013, an era range containing long
stretches of 80--90\% unanimity that no modern model enjoys; on the
modern, 5--4-heavy window, the honest point of comparison is
closer to our setting, and the 63.7\% bar is what an
ideology-only adversary achieves when unanimity is gone.

\subsection{Why the null makes the text conditions more
interesting, not less}

Had a structured challenger cleared the bar, the persona experiment
would have been partially redundant: some of the textual signal
would already be recoverable from codes. The null result does the
opposite---it concentrates the entire residual question in the
text. If prediction above 63.7\% is possible at all on this window,
it must come from information that is written down but not coded:
the argument's framing, the questions the justices asked at oral
argument, the linguistic signature of each opinion. The persona
condition (B) tests the strongest version of that hypothesis ---
that the signature is \emph{individual} and \emph{prospective} ---
and the stylometry lineage~\citep{mosteller1964,juola2006} gives it
a principled footing: style is known to be persistent and hard to
fake, so a fine-tuned adapter is a reasonable instrument for
measuring whether style and substance are entangled.

\subsection{Calibrating expectations for the Final Test}

The sealed fifty are 5--4 decisions, which the folklore of court
watching regards as the least predictable class; and the bar was
computed on the full evaluation window, not on 5--4s specifically.
Both facts cut against the persona condition clearing B4 on the
holdout, and the honest prior is that the decisive test may return
``no difference''. We stress that this is a designed outcome, not a
failure state: a tight null on a sealed, pre-registered holdout
with per-justice adapters and audited anti-memorization checks
would be the cleanest public evidence yet that judicial
subjectivity, beyond ideology, is not extractible from the written
record---a claim that would stand regardless of how much compute
a future replicant throws at it.
